\section{Introduction}\label{sec:introduction}

The triangular Hilbert transform couples translations in two different
coordinate directions. For complex functions $F,G$ on $\R^2$ and
$0<\varepsilon<R<\infty$, write
\[
 B_{\varepsilon,R}(F,G)(x,y)
 =\int_{\varepsilon<|t|<R}F(x+t,y)G(x,y+t)\,\frac{\dd t}{t},
 \qquad
 B_*(F,G)=\sup_{0<\varepsilon<R<\infty}|B_{\varepsilon,R}(F,G)|.
\]
For $L^3$ inputs, the finite integrals are absolutely convergent outside
one common null set, and the supremum is measurable; the details are
given in Lemma~\ref{lem:finite-truncations}.

\begin{theorem}\label{thm:main}
There is an absolute constant $C$ such that, for every pair of complex
functions $F,G\in L^3(\R^2)$,
\[
 \|B_*(F,G)\|_{L^{3/2}(\R^2)}
 \le C\|F\|_{L^3(\R^2)}\|G\|_{L^3(\R^2)}.
\]
The supremum is over both finite hard truncation endpoints.
\end{theorem}

The estimate controls a truncation chosen separately at each spatial
point. This is stronger than a uniform bound for the norm of each
individual truncation: the same $L^{3/2}$ majorant must control every
choice of the two endpoints. It also gives the joint principal value
\[
 B(F,G)=\lim_{\substack{\varepsilon\downarrow0\\R\uparrow\infty}}
 B_{\varepsilon,R}(F,G)
\]
almost everywhere and in $L^{3/2}$, with the same norm bound
(Theorem~\ref{thm:bilinear-pv}).

Pairing the bilinear output with a third $L^3$ function leads to the
scalar triangular form. In its symmetric coordinates this is
\[
 \Lambda_{\varepsilon,R}(G_0,G_1,G_2)
 =\iiint_{\varepsilon<|a+b+c|<R}
 G_0(a,b)G_1(b,c)G_2(c,a)
 \,\frac{\dd a\,\dd b\,\dd c}{a+b+c}.
\]
The three factors have equal exponent $3$, which is the symmetric point
of the scalar H\"older relation $1/p_0+1/p_1+1/p_2=1$.
We obtain an endpoint-uniform bound by $C\prod_v\|G_v\|_3$ and a
joint scalar principal value for every complex $L^3$ triple.
Proposition~\ref{prop:flat-equivalence} gives the explicit
determinant-one changes of variables and proves that the flat and
simplex scalar estimates have the same optimal constant.
In particular, the result resolves positively the symmetric triangular
Hilbert transform problem stated by Thiele
in~\cite[Section~9, Problem~13]{GrafakosEtAl2017}.

\subsection{Previous work}

The flat operator appears in Demeter and Thiele's discussion of
bilinear averages for two commuting transformations~\cite[Section~6]{DemeterThiele2010}.
The three inputs of the scalar form depend on the three pairs of
variables, making it a basic example of an entangled singular integral.
Kova\v{c}, Thiele and Zorin-Kranich proved estimates for a Walsh model
with one structurally restricted input~\cite[Theorem~1.7]{KovacThieleZorinKranich2015}.

Tao proved sublogarithmic cancellation for one-dimensional multilinear
Hilbert transforms~\cite[Theorem~1.2]{Tao2016Cancellation}.
Zorin-Kranich extended this cancellation to simplex Hilbert
transforms~\cite[Theorem~1.3]{ZorinKranich2017}. Square-function estimates of Durcik,
Kova\v{c}, \v{S}kreb and Thiele gave square-root cancellation for the
triangular form across finitely many smooth scales~\cite[Corollary~5]{DurcikKovacSkrebThiele2019}.
Durcik, Kova\v{c} and Thiele proved power-type cancellation with hard
truncations~\cite[Theorem~1]{DurcikKovacThiele2019}; at the symmetric
tuple the bound grows like $\sqrt{\log(R/\varepsilon)}$, by permutation
and interpolation of their $(4,4,2)$ estimates.
These finite-scale bounds retain a dependence on the scale range.

Bounds for averages over direction parameters were obtained by
Durcik and Roos~\cite[Theorem~1]{DurcikRoos2021} and, for rough angular
weights, by Lin and Slav\'ikov\'a~\cite[Theorem~1.1]{LinSlavikova2026}.
For curved translations, Christ, Durcik and Roos proved bounds along
$(t,t^2)$~\cite[Theorem~1]{ChristDurcikRoos2021}. Hsu and Lin treat
more general curves under geometric hypotheses excluding the flat curve
$(t,t)$~\cite[Theorem~1.37]{HsuLin2026}. The curvature in these results
provides an additional source of cancellation that is unavailable for
the flat curve.

Our argument uses a nonnegative matrix energy and Gaussian heat flow.
Energy telescoping for entangled forms appears in Kova\v{c}'s
twisted-paraproduct argument~\cite[Lemma~2]{Kovac2012Twisted} and his
Bellman-function framework~\cite[Section~2]{Kovac2011}. Continuous
Gaussian telescoping appears in Durcik~\cite[Lemma~3 and Section~3]{Durcik2015}.
These methods concern bipartite forms; Kova\v{c} discusses the obstacle
posed by the triangular cycle in~\cite[Section~6.2]{Kovac2011}.
The cyclic trace representation itself also occurs
in~\cite[Equation~(1.5)]{KovacThieleZorinKranich2015}.
We prove a dimension-independent mixed trace inequality and compare
the energy derivatives at neighboring vertices. The key additional
estimate controls the total change in energy when individual matrix
entries are switched on and off at independently chosen scales.
This permits the pointwise maximum without a factor depending on
the number of scales or the matrix dimensions.

\subsection{Proof overview}

The change of variables $x=-a$, $y=-b$, $t=a+b+c$ writes the bilinear
output in terms of
\[
 G_1(b,c)=F(b+c,-b),\qquad G_2(c,a)=G(-a,c+a).
\]
Both pair maps preserve $L^3$ norms. We first sample these inputs,
together with a dual input $G_0(a,b)$, on a finite mesh. Multiplying
entries by Gaussian weights gives maps $W_v:H_{v+1}\to H_v$ between
finite-dimensional complex Hilbert spaces, indexed by
$v\in\mathbb Z/3\mathbb Z$.
The three vertex spaces correspond to the sampled $a$-, $b$- and
$c$-coordinates. Thus the two indices of the dual array record the
output position $(x,y)=(-a,-b)$.
At a vertex, combine its two incident maps into
$R_v=[\,W_v\quad W_{v-1}^*\,]$ and use the energy
$\tr((R_v^*R_v)^{3/2})$.

Sections~\ref{sec:matrix} and~\ref{sec:trace} prove the matrix
inequalities. Section~\ref{sec:heat} then integrates the sum of the
three energies over Gaussian centers satisfying $p_0+p_1+p_2=0$.
For fixed arrays, its decrease as the variance grows controls the quadratic
costs needed to estimate a cyclic trace with three odd Gaussian windows.

To select different truncations at different output points, the dual
array entry at $(i,j)$ must be present only between its own two scales.
The energy is therefore decreasing between finitely many switching
scales and may jump at those scales. Section~\ref{sec:jumps} bounds
the sum of the absolute jumps by the discrete $L^3$ masses of the three
arrays. The underlying matrix fact is that the derivative with respect
to an entry is controlled by the product of its row and column norms.
After Gaussian integration, one-dimensional maximal functions bound
these products uniformly in the switching scales.

Section~\ref{sec:maximal} uses this bound to linearize any finite
maximum of smooth truncations. Mesh limits and rational exhaustion
give the full smooth maximum. A pointwise comparison of its kernel
with the hard Hilbert kernel completes the estimate for compactly
supported smooth inputs. Finally, Section~\ref{sec:consequences}
uses density and a vanishing tail-oscillation argument to prove
Theorem~\ref{thm:main}, the joint principal values, and the scalar
equivalence.
